\section{Thiết kế cơ sở dữ liệu chi tiết (Physical Database Design)}
\label{sec:database_design}

Cơ sở dữ liệu trung tâm của hệ thống (\code{clinical_hospital.db}) được thiết kế theo mô hình quan hệ đã chuẩn hóa đến dạng chuẩn 3 (3NF), lưu trữ cục bộ và tối ưu hóa ở mức nhị phân.

\subsection{Cấu hình hiệu năng cao SQLite WAL Mode}
Hệ thống kích hoạt các thiết lập pragma tối ưu tại thời điểm khởi tạo kết nối:
\begin{itemize}[leftmargin=1.5em]
    \item \code{PRAGMA journal_mode = WAL;} Cho phép cơ chế Write-Ahead Logging; các thao tác đọc và ghi diễn ra đồng thời mà không khóa lẫn nhau.
    \item \code{PRAGMA synchronous = NORMAL;} Giảm thiểu tần suất ghi đĩa đồng bộ, tăng tốc độ xử lý I/O đĩa lên gấp 3 lần.
    \item \code{PRAGMA cache_size = -64000;} Cấp phát 64MB bộ nhớ RAM làm bộ đệm trang (Page Cache) cho CSDL.
\end{itemize}

\subsection{Từ điển dữ liệu chi tiết các bảng (Data Dictionary)}

\begin{table}[htbp]
\centering
\small
\begin{tabularx}{\textwidth}{L{3.5cm} L{2.2cm} C{1.2cm} C{1.8cm} Y}
\toprule
\textbf{Tên cột} & \textbf{Kiểu dữ liệu} & \textbf{Khóa} & \textbf{Ràng buộc} & \textbf{Mô tả nghiệp vụ thực tế} \\
\midrule
\code{id} & INTEGER & PK & NOT NULL & Khóa chính định danh tự tăng. \\
\code{registration_no} & VARCHAR(50) & & NOT NULL & Số đăng ký lưu hành do DAV cấp (VD: VD-25489-16). \\
\code{drug_name} & VARCHAR(255) & & NOT NULL & Tên biệt dược thương mại của thuốc. \\
\code{active_ingredient} & VARCHAR(255) & & NOT NULL & Tên hoạt chất quốc tế INN chuẩn hóa. \\
\code{content} & VARCHAR(100) & & & Nồng độ hoặc hàm lượng thuốc (VD: 500mg, 1g). \\
\code{dosage_form} & VARCHAR(100) & & & Dạng bào chế (viên nén, dung dịch tiêm). \\
\code{manufacturer} & VARCHAR(255) & & & Tên công ty/nhà máy sản xuất dược phẩm. \\
\code{country} & VARCHAR(100) & & & Quốc gia sản xuất (Việt Nam, Pháp, Ấn Độ...). \\
\bottomrule
\end{tabularx}
\caption{Từ điển dữ liệu bảng \code{dav_drugs} (54.885 thuốc Cục Quản lý Dược)}
\end{table}

\begin{table}[htbp]
\centering
\small
\begin{tabularx}{\textwidth}{L{3.2cm} L{2.2cm} C{1.2cm} C{1.8cm} Y}
\toprule
\textbf{Tên cột} & \textbf{Kiểu dữ liệu} & \textbf{Khóa} & \textbf{Ràng buộc} & \textbf{Mô tả nghiệp vụ thực tế} \\
\midrule
\code{id} & INTEGER & PK & NOT NULL & Khóa chính định danh tự tăng. \\
\code{drug_a} & VARCHAR(100) & & NOT NULL & Hoạt chất INN thứ nhất (đã chuẩn hóa chữ thường). \\
\code{drug_b} & VARCHAR(100) & & NOT NULL & Hoạt chất INN thứ hai (đã chuẩn hóa chữ thường). \\
\code{severity_level} & VARCHAR(20) & & NOT NULL & Mức độ nghiêm trọng: CONTRAINDICATED (Cấm), MAJOR. \\
\code{mechanism} & TEXT & & NOT NULL & Cơ chế dược lý sinh học dẫn đến tương tác. \\
\code{clinical_effect} & TEXT & & NOT NULL & Hậu quả lâm sàng (tử vong, xuất huyết, suy thận cấp). \\
\code{legal_reference} & VARCHAR(100) & & NOT NULL & Căn cứ pháp lý (Quyết định 5948/QĐ-BYT). \\
\bottomrule
\end{tabularx}
\caption{Từ điển dữ liệu bảng \code{ddi_rules} (Tương tác chống chỉ định)}
\end{table}

\begin{table}[htbp]
\centering
\small
\begin{tabularx}{\textwidth}{L{3.2cm} L{2.2cm} C{1.2cm} C{1.8cm} Y}
\toprule
\textbf{Tên cột} & \textbf{Kiểu dữ liệu} & \textbf{Khóa} & \textbf{Ràng buộc} & \textbf{Mô tả nghiệp vụ thực tế} \\
\midrule
\code{id} & VARCHAR(50) & PK & NOT NULL & Mã bệnh nhân (Hospital Patient ID). \\
\code{full_name} & VARCHAR(150) & & NOT NULL & Họ và tên đầy đủ của người bệnh. \\
\code{dob} & DATE & & NOT NULL & Ngày tháng năm sinh. \\
\code{gender} & VARCHAR(10) & & NOT NULL & Giới tính (Nam, Nữ). \\
\code{bhyt_code} & VARCHAR(15) & & & Mã số thẻ BHYT 15 ký tự theo Quyết định 1666. \\
\code{bed_number} & VARCHAR(20) & & & Số giường bệnh nội trú đang nằm điều trị. \\
\code{serum_cr} & REAL & & & Chỉ số Creatinine huyết thanh gần nhất (mg/dL). \\
\code{egfr} & REAL & & & Mức lọc cầu thận tính toán theo Cockcroft-Gault. \\
\bottomrule
\end{tabularx}
\caption{Từ điển dữ liệu bảng \code{patients} (Quản lý người bệnh nội trú)}
\end{table}

\begin{table}[htbp]
\centering
\small
\begin{tabularx}{\textwidth}{L{3.2cm} L{2.2cm} C{1.2cm} C{1.8cm} Y}
\toprule
\textbf{Tên cột} & \textbf{Kiểu dữ liệu} & \textbf{Khóa} & \textbf{Ràng buộc} & \textbf{Mô tả nghiệp vụ thực tế} \\
\midrule
\code{id} & INTEGER & PK & NOT NULL & Khóa chính tự tăng của bản ghi nhật ký. \\
\code{user_id} & VARCHAR(50) & & NOT NULL & Mã tài khoản nhân viên y tế thực hiện thao tác. \\
\code{action_type} & VARCHAR(50) & & NOT NULL & Loại hành động (PRESCRIBE, DDI_ALERT, OVERRIDE_ALERT). \\
\code{patient_id} & VARCHAR(50) & FK & NOT NULL & Mã bệnh nhân liên quan đến thao tác. \\
\code{payload} & TEXT & & NOT NULL & Dữ liệu chi tiết của sự kiện dưới định dạng JSON. \\
\code{timestamp} & DATETIME & & NOT NULL & Thời điểm chính xác xảy ra sự kiện theo chuẩn ISO 8601. \\
\bottomrule
\end{tabularx}
\caption{Từ điển dữ liệu bảng \code{audit_logs} (Nhật ký kiểm toán lâm sàng)}
\end{table}

\begin{table}[htbp]
\centering
\small
\begin{tabularx}{\textwidth}{L{3.5cm} L{2.2cm} C{1.2cm} C{1.8cm} Y}
\toprule
\textbf{Tên cột} & \textbf{Kiểu dữ liệu} & \textbf{Khóa} & \textbf{Ràng buộc} & \textbf{Mô tả nghiệp vụ thực tế} \\
\midrule
\code{id} & VARCHAR(50) & PK & NOT NULL & Mã định danh nhân viên y tế (Mã NV). \\
\code{name} & VARCHAR(150) & & NOT NULL & Họ tên nhân viên kèm học hàm/học vị. \\
\code{role} & VARCHAR(50) & & NOT NULL & Vai trò phân quyền (doctor, nurse, pharmacist, admin). \\
\code{department} & VARCHAR(100) & & NOT NULL & Khoa lâm sàng công tác (Nội tổng hợp, Cấp cứu...). \\
\code{practice_cert} & VARCHAR(50) & & & Số chứng chỉ hành nghề khám chữa bệnh (CCHN). \\
\code{username} & VARCHAR(100) & & UNIQUE & Tên đăng nhập Web (@username). \\
\code{hashed_password} & VARCHAR(255) & & & Mật khẩu băm chuẩn quân sự Argon2id. \\
\code{must_change_pwd} & BOOLEAN & & DEFAULT 0 & Cờ bắt buộc đổi mật khẩu ở lần đăng nhập đầu. \\
\code{rfid_card_id} & VARCHAR(50) & & & Mã định danh thẻ từ RFID quẹt ca trực. \\
\code{pin_code} & VARCHAR(10) & & & Mã PIN ca trực 4 số thao tác tại giường bệnh. \\
\code{is_active} & BOOLEAN & & DEFAULT 1 & Trạng thái tài khoản (1: Đang công tác, 0: Tạm ngưng). \\
\bottomrule
\end{tabularx}
\caption{Từ điển dữ liệu bảng \code{users} (Quản lý hồ sơ nhân sự và phân quyền bảo mật)}
\end{table}

\begin{table}[htbp]
\centering
\small
\begin{tabularx}{\textwidth}{L{3.5cm} L{2.2cm} C{1.2cm} C{1.8cm} Y}
\toprule
\textbf{Tên cột} & \textbf{Kiểu dữ liệu} & \textbf{Khóa} & \textbf{Ràng buộc} & \textbf{Mô tả nghiệp vụ thực tế} \\
\midrule
\code{session_id} & VARCHAR(64) & PK & NOT NULL & Mã phiên hội chẩn AI định danh UUID v4. \\
\code{user_id} & VARCHAR(50) & FK & NOT NULL & Mã tài khoản nhân sự sở hữu phiên hội chẩn. \\
\code{patient_id} & VARCHAR(50) & & & Mã bệnh nhân liên quan đến ca hội chẩn (nếu có). \\
\code{title} & VARCHAR(255) & & NOT NULL & Tiêu đề tóm tắt ca tư vấn lâm sàng. \\
\code{created_at} & DATETIME & & NOT NULL & Thời điểm khởi tạo phiên trên đám mây CSDL. \\
\code{updated_at} & DATETIME & & NOT NULL & Thời điểm tin nhắn cuối cùng được trao đổi. \\
\bottomrule
\end{tabularx}
\caption{Từ điển dữ liệu bảng \code{chat_sessions} (Lưu trữ phiên hội chẩn đám mây theo tài khoản)}
\end{table}

\begin{table}[htbp]
\centering
\small
\begin{tabularx}{\textwidth}{L{3.5cm} L{2.2cm} C{1.2cm} C{1.8cm} Y}
\toprule
\textbf{Tên cột} & \textbf{Kiểu dữ liệu} & \textbf{Khóa} & \textbf{Ràng buộc} & \textbf{Mô tả nghiệp vụ thực tế} \\
\midrule
\code{id} & INTEGER & PK & NOT NULL & Khóa chính tự tăng của tin nhắn. \\
\code{session_id} & VARCHAR(64) & FK & NOT NULL & Mã phiên tham chiếu đến bảng \code{chat_sessions}. \\
\code{role} & VARCHAR(20) & & NOT NULL & Vai trò người gửi (\code{user} hoặc \code{assistant}). \\
\code{content} & TEXT & & NOT NULL & Nội dung tin nhắn hội chẩn hoặc câu trả lời AI. \\
\code{citations} & TEXT & & & Dữ liệu JSON chứa con trỏ trích dẫn số trang Dược thư. \\
\code{created_at} & DATETIME & & NOT NULL & Thời điểm gửi tin nhắn theo chuẩn ISO 8601. \\
\bottomrule
\end{tabularx}
\caption{Từ điển dữ liệu bảng \code{chat_messages} (Chi tiết tin nhắn hội chẩn lâm sàng)}
\end{table}

\begin{table}[htbp]
\centering
\small
\begin{tabularx}{\textwidth}{L{3.5cm} L{2.2cm} C{1.2cm} C{1.8cm} Y}
\toprule
\textbf{Tên cột} & \textbf{Kiểu dữ liệu} & \textbf{Khóa} & \textbf{Ràng buộc} & \textbf{Mô tả nghiệp vụ thực tế} \\
\midrule
\code{id} & VARCHAR(50) & PK & NOT NULL & Mã định danh văn bản pháp quy nạp vào hệ thống. \\
\code{title} & VARCHAR(255) & & NOT NULL & Tên văn bản pháp quy/hướng dẫn điều trị (PDF). \\
\code{file_hash} & VARCHAR(64) & & NOT NULL & Mã băm SHA-256 chống trùng lặp dữ liệu nạp. \\
\code{chunk_count} & INTEGER & & NOT NULL & Tổng số phân đoạn ngữ nghĩa (Chunks) sau khi cắt. \\
\code{status} & VARCHAR(30) & & NOT NULL & Trạng thái xử lý (EMBEDDED, READY, ARCHIVED). \\
\code{created_at} & DATETIME & & NOT NULL & Thời điểm nạp tài liệu vào CSDL RAG. \\
\bottomrule
\end{tabularx}
\caption{Từ điển dữ liệu bảng \code{ingested_documents} (Quản lý tài liệu pháp quy nạp động)}
\end{table}
